\section{SECCIÓN 12. GESTIÓN DE RIESGOS}
La gestión de riesgos tiene como objetivo identificar los posibles eventos que pueden afectar la calidad del software o el desarrollo del proceso de pruebas, estableciendo acciones preventivas y correctivas para minimizar su impacto.

\subsection*{Riesgos del Producto}
\begin{longtable}{|p{1.5cm}|p{3.5cm}|p{2.5cm}|p{2cm}|p{2cm}|p{4cm}|}
\hline
\textbf{ID} & \textbf{Riesgo} & \textbf{Probabilidad} & \textbf{Impacto} & \textbf{Nivel} & \textbf{Estrategia de Mitigación} \\
\hline
RP-01 & Acceso no autorizado por vulnerabilidades en JWT o RBAC & Media & Muy Alto & Crítico & Validar autenticación, autorización y expiración de tokens mediante pruebas de seguridad. \\
\hline
RP-02 & Ataques XSS mediante contenido enviado por el usuario & Media & Alto & Alto & Aplicar sanitización con OWASP Java HTML Sanitizer y realizar pruebas de penetración. \\
\hline
RP-03 & Ataques CSRF sobre operaciones protegidas & Baja & Alto & Medio & Implementar y verificar el funcionamiento del filtro CSRF con Double Submit Cookie. \\
\hline
RP-04 & Pérdida de registros de auditoría & Baja & Alto & Medio & Registrar todas las operaciones críticas mediante AuditLog y respaldar la base de datos periódicamente. \\
\hline
RP-05 & Fallos en el almacenamiento seguro de imágenes & Baja & Alto & Medio & Utilizar nombres UUID, directorios Safe y Quarantine y validar permisos de acceso. \\
\hline
RP-06 & Errores en la validación de archivos (tipo o tamaño) & Media & Medio & Medio & Verificar tamaño máximo de 50 MB, tipos MIME permitidos y reglas de validación antes del upload. \\
\hline
\end{longtable}

\subsection*{Riesgos del Proyecto de Pruebas}
\begin{longtable}{|p{1.5cm}|p{3.5cm}|p{2.5cm}|p{2cm}|p{2cm}|p{4cm}|}
\hline
\textbf{ID} & \textbf{Riesgo} & \textbf{Probabilidad} & \textbf{Impacto} & \textbf{Nivel} & \textbf{Estrategia de Mitigación} \\
\hline
RPR-01 & Retrasos en el cronograma de pruebas & Media & Alto & Alto & Planificar las actividades por Sprint y realizar reuniones de seguimiento semanales. \\
\hline
RPR-02 & Casos de prueba incompletos & Media & Alto & Alto & Elaborar una matriz de trazabilidad entre requisitos y casos de prueba antes de iniciar la ejecución. \\
\hline
RPR-03 & Cambios frecuentes en el código durante las pruebas & Alta & Medio & Alto & Utilizar GitHub con ramas de desarrollo y control de versiones para gestionar cambios. \\
\hline
RPR-04 & Indisponibilidad de PostgreSQL o ClamAV & Baja & Alto & Medio & Mantener un entorno de pruebas estable y verificar los servicios antes de ejecutar las pruebas. \\
\hline
RPR-05 & Baja cobertura de pruebas automatizadas & Media & Medio & Medio & Implementar pruebas unitarias con JUnit y pruebas de API mediante Postman. \\
\hline
RPR-06 & Defectos no registrados correctamente & Baja & Medio & Bajo & Mantener un registro centralizado de incidencias con evidencia, prioridad y estado. \\
\hline
\end{longtable}

\subsection*{Plan de Respuesta a los Riesgos}
Para minimizar el impacto de los riesgos identificados durante el proceso de aseguramiento de la calidad, se implementarán las siguientes acciones utilizando las herramientas definidas en el proyecto:
\begin{itemize}
    \item Ejecutar análisis estáticos del código mediante SonarQube para identificar vulnerabilidades, code smells, errores potenciales y problemas de mantenibilidad antes de iniciar las pruebas dinámicas.
    \item Medir la cobertura de las pruebas utilizando JaCoCo, con el fin de identificar componentes del sistema que no han sido suficientemente validados y mejorar la confiabilidad del software.
    \item Analizar las dependencias del proyecto mediante npm audit para detectar bibliotecas vulnerables o desactualizadas y aplicar las actualizaciones o correcciones necesarias.
    \item Ejecutar pruebas unitarias con JUnit para validar el correcto funcionamiento de los componentes críticos del backend, asegurando que la lógica de negocio opere de acuerdo con los requisitos establecidos.
    \item Realizar pruebas de seguridad utilizando OWASP ZAP, con el objetivo de identificar vulnerabilidades en la aplicación web, como problemas de autenticación, control de acceso, inyección de código o configuraciones inseguras.
    \item Evaluar el rendimiento, la accesibilidad y las buenas prácticas del frontend mediante Lighthouse, identificando oportunidades de mejora que contribuyan a una mejor experiencia de usuario.
    \item Documentar todos los hallazgos, defectos y vulnerabilidades detectadas durante el proceso de pruebas, verificando posteriormente las correcciones mediante pruebas de regresión para asegurar que los cambios implementados no introduzcan nuevos errores.
\end{itemize}